\DocumentMetadata{lang=en-US,pdfstandard=A-2b,uncompress}
\documentclass[12pt]{article}
\usepackage[utf8]{inputenc}
\usepackage[margin=1in]{geometry}
\usepackage{amsmath,amssymb}
\usepackage{enumitem}
\usepackage{hyperref}
\usepackage[tagged, highstructure]{accessibility}
\usepackage{graphicx}
\usepackage{fitch}

\title{Quiz 2: DNF and Soundness}
\hypersetup{pdftitle={Quiz 2: DNF and Soundness}}
\author{}
\date{}

\newcommand{\forces}{\models}
\newcommand{\meta}[1]{\mathcal{#1}}

\begin{document}
\maketitle

\textbf{Instructions:} Show your reasoning. Write clearly enough that another student could follow your argument.

\begin{enumerate}
    \item \textbf{DNF via the truth-table method}
    \begin{enumerate}[label=(\alph*)]
        \item Use the general truth-table method to put the sentence
        \[
        (P \rightarrow R)\leftrightarrow R
        \]
        into \textbf{disjunctive normal form} (DNF).

        \item Explain why your solution works.
        \begin{itemize}
            \item How does each disjunct relate to rows of the original sentence's truth table?
            \item How does the number of disjuncts relate to the truth table?
        \end{itemize}
    \end{enumerate}
    \textit{Reminder (DNF method):} use the rows where the original sentence is \textbf{true}; for each such row build a conjunction true only on that row; then disjoin those conjunctions.

    \item \textbf{Reading a proof using \(\delta_i\), \(\phi_i\), and shininess}

    Consider the following TFL Fitch proof.

    \begin{center}
    \[
    \begin{nd}
    \open
      \hypo{1}{P \rightarrow Q}
      \open
        \hypo{2}{\neg Q}
        \open
          \hypo{3}{P}
          \have{4}{Q}\by{$\rightarrow$E}{1,3}
          \have{5}{\bot}\ne{2,4}
        \close
        \have{6}{\neg P}\ni{3-5}
      \close
        \have{7}{\neg Q \rightarrow \neg P}\by{$\rightarrow$I}{2-6}
    \close
    \have{8}{(P \rightarrow Q)\rightarrow(\neg Q\rightarrow\neg P)}\by{$\rightarrow$I}{1-7}
    \end{nd}
    \]
    \end{center}

    \begin{enumerate}[label=(\alph*)]
        \item What are \(\delta_6\), \(\phi_6\), \(\delta_8\), and \(\phi_8\) for this proof?
        \item What does it mean for line \(n\) to be \textbf{shiny}?
        \item What does it mean for a rule to be \textbf{rule-sound}, in terms of shininess?
    \end{enumerate}

    \item \textbf{Fill in missing justifications: show IP is rule-sound}

    Assume all lines \(<n\) are shiny, and line \(n\) is justified by IP from subproof \(i\text{-}j\), where:
    \[
    \phi_i=\neg A,\quad \phi_j=\bot,\quad \phi_n=A.
    \]

    Schematic proof shape for IP:
    \begin{center}
      \alt{Indirect proof schema: assume not A, derive contradiction, conclude A by IP.}
      \includegraphics[width=0.42\linewidth]{proof-images/Complete_Proof_Rules/proof_15_indirect.png}
    \end{center}

    Fill in the missing justifications.
    \begin{enumerate}[label=\arabic*.]
        \item There is some \(\delta'\) such that \(\delta' \subseteq \delta_n\) and \(\delta_j=\delta' \cup \{\neg A\}\). \hfill [Justification: \rule{2.8in}{0.4pt}]
        \item Since \(j<n\), line \(j\) is shiny, so \(\delta_j \forces \bot\). \hfill [Justification: \rule{2.8in}{0.4pt}]
        \item Let \(v\) be any valuation satisfying \(\delta_n\). We must show \(v \forces A\).
        \item Assume for contradiction that \(v \not\forces A\). Then \(v \forces \neg A\). \hfill [Justification: \rule{2.8in}{0.4pt}]
        \item \hspace*{1.5em}So \(v \forces \delta_j\). \hfill [Justification: \rule{2.8in}{0.4pt}]
        \item \hspace*{1.5em}Hence \(v \forces \bot\), contradiction. \hfill [Justification: \rule{2.8in}{0.4pt}]
        \item Therefore \(v \forces A\). Since \(v\) was arbitrary, \(\delta_n \forces A\), so line \(n\) is shiny.
    \end{enumerate}

    \item \textbf{Prove that \( \vee I \) is rule-sound.}

    Write a full proof in the same style used in class:
    \begin{itemize}
        \item Start with the inductive hypothesis setup.
        \item State the relevant \(\delta_i \subseteq \delta_n\)-type fact(s).
        \item Use an arbitrary valuation satisfying \(\delta_n\).
        \item Conclude line \(n\) is shiny.
    \end{itemize}

    \item \textbf{Extra credit:} Prove that \( \vee E \) is rule-sound.
\end{enumerate}

\end{document}
